\documentclass[11pt,a4paper]{article}
\usepackage[margin=2.4cm]{geometry}
\usepackage{amsmath,amssymb,booktabs,hyperref}
\newcommand{\dt}{\Delta t}
\newcommand{\avg}[1]{\{\!\{#1\}\!\}}
\newcommand{\jump}[1]{[\![#1]\!]}
\newcommand{\code}[1]{\texttt{#1}}
\title{The frozen split of the finite-volume Euler equations\\on the C-grid}
\author{MultiRate notes}
\date{\today}
\begin{document}
\maketitle

\begin{abstract}
For split-explicit (MIS) time integration with a \emph{frozen} right-hand side, $N(q)=F(q,q)$, the
first argument $y$ is frozen at the stage values and the second $z$ evolves in the fast loop. This note
describes the split of the finite-volume discretization on the Arakawa C-grid
(\code{src/FV/frozen.jl}, flat grids): the potential temperature at the faces and, optionally, the
derivative $\mathrm dp/\mathrm d\Theta$ at the faces are frozen, the momentum advection is the slow part. Two FV
schemes are provided, with the pressure gradient $\nabla p$ or in product form $(\mathrm dp/\mathrm d\Theta)\nabla\Theta$; for the
product form the analytic $\mathrm dp/\mathrm d\Theta$ is frozen, the fast part is linear in $z$ and its vertically implicit (HEVI) step is a
tridiagonal solve per column. A well-balanced variant without background subtraction, following
Artiano et al., is given, and the consequences for the stability of MIS methods are summarized.
\end{abstract}

\section{Equations}
Variables $\rho$, $U=\rho u$, $W=\rho w$, $\Theta=\rho\theta$; the pressure is a function of $\Theta$ alone,
\[
p=P(\Theta)=p_0\Big(\frac{R\,\Theta}{p_0}\Big)^{\gamma},\qquad P'(\Theta)=\frac{\gamma p}{\Theta}=\frac{c^2}{\theta}.
\]
\begin{align*}
\partial_t\rho+\partial_xU+\partial_zW&=0,\\
\partial_t\Theta+\partial_x(\theta U)+\partial_z(\theta W)&=0,\\
\partial_tU+\partial_x(Uu)+\partial_z(Wu)+\partial_xp&=0,\\
\partial_tW+\partial_x(Uw)+\partial_z(Ww)+\partial_zp+\rho g&=0,
\end{align*}
plus a Rayleigh sponge $-\alpha(z)(q-q_{\rm target})$ near the top (and optionally at the sides). No
background state is subtracted (the optional \code{FV\_HYDRO\_REFERENCE} subtracts $p_{\rm ref}$,
$\rho_{\rm ref}$ in the momentum equations; it is off by default, and the frozen split rejects such models).

\section{Grid and full operator $N$}
Flat C-grid, $n_x\times n_z$ cells of size $\Delta x\times\Delta z$, periodic in $x$, walls at the bottom and
the top. $\rho_{i,k}$ and $\Theta_{i,k}$ at the cell centres, $U_{i-1/2,k}$ at the $x$-faces, $W_{i,k-1/2}$ at the
interior $z$-faces ($W=0$ at the walls). Velocities at the faces divide by the arithmetic mean of the two
cells, $u_{i-1/2,k}=U_{i-1/2,k}/\avg\rho_x$, $w_{i,k-1/2}=W_{i,k-1/2}/\avg\rho_z$; $\theta=\Theta/\rho$ at the centres.
Differences $\delta_x a_{i}=(a_{i+1/2}-a_{i-1/2})/\Delta x$ and likewise in $z$.

\paragraph{Face values.} At every face $i+\tfrac12$ a left and a right value are reconstructed from the
cell values by the third-order ($\kappa=\tfrac13$) formula,
\[
a^L_{i+1/2}=\frac{-a_{i-1}+5a_i+2a_{i+1}}6,\qquad
a^R_{i+1/2}=\frac{2a_i+5a_{i+1}-a_{i+2}}6,
\]
and the upwind one is used,
\[
a_{i+1/2}=\begin{cases}a^L_{i+1/2},&v\ge0,\\ a^R_{i+1/2},&v<0,\end{cases}
\]
$v$ the transporting mass flux at the face (\code{recon\_left}, \code{recon\_right}, \code{interp3} in
\code{src/FV/advection.jl}). (This is the same value as the third-order upwind-biased flux of WRF, written there
as a centred fourth-order mean plus a sign-dependent correction; the reconstruction form makes the left and
right states explicit.) In $z$ at the faces next to the walls, where one of the two states lacks its
stencil, the centred mean $\tfrac12(a_i+a_{i+1})$ is used.

\paragraph{The operator.}
\begin{align*}
\partial_t\rho_{i,k}&=-\delta_xU-\delta_zW,\\
\partial_t\Theta_{i,k}&=-\delta_x(\tilde\theta_xU)-\delta_z(\tilde\theta_zW),\qquad
\tilde\theta_x,\ \tilde\theta_z\ \text{the upwind face values of }\theta\ \text{for }U,\ W,\\
\partial_tU_{i-1/2,k}&=-A_U-\frac{p_{i,k}-p_{i-1,k}}{\Delta x},\\
\partial_tW_{i,k-1/2}&=-A_W-\frac{p_{i,k}-p_{i,k-1}}{\Delta z}-g\,\frac{\rho_{i,k-1}+\rho_{i,k}}2,
\end{align*}
$p_{i,k}=P(\Theta_{i,k})$ (scheme \code{:gradp}, the default of \code{FV\_PRESSURE\_FORM}). The product-form
scheme \code{:product} replaces the pressure differences by
\[
p_{i,k}-p_{i-1,k}\ \to\ P'_{i-1/2,k}\,(\Theta_{i,k}-\Theta_{i-1,k}),\qquad
P'_{i-1/2,k}=\tfrac12\big(P'(\Theta_{i-1,k})+P'(\Theta_{i,k})\big),\quad P'(\Theta)=\frac{\gamma P(\Theta)}{\Theta},
\]
and likewise in $z$: the derivative is computed analytically in the cells and averaged to the face. It is
defined for all neighbouring values (a secant $\jump p/\jump\Theta$ is not, for equal values), and its pressure
difference differs from $\jump p$ by $O(\jump\Theta^3)$. Both schemes use the full $p$ and $\rho$, no background
state. The momentum advection $A_U$, $A_W$ is in flux form with the transporting mass fluxes
averaged to the momentum control volumes (e.g.\ $\bar U_{i,k}=\tfrac12(U_{i-1/2,k}+U_{i+1/2,k})$) and the
transported velocities as upwind face values (the same reconstruction); $A_W$ uses the centred mean next to the walls.

\section{The frozen split}
\paragraph{Frozen coefficients} $C(y)$, computed once per stage value $Y_j$ (\code{frozen\_data!}):
\begin{itemize}
\item $\tilde\theta_x(y)$, $\tilde\theta_z(y)$: the upwinded face values of $\theta$ of the full operator, the upwind
  direction from $U$, $W$ of $y$;
\item for the scheme \code{:product} the face derivatives $P'_f(y)$ (the means of the cell values $\gamma p/\Theta$ of
  $y$, as in $N$).
\end{itemize}

\paragraph{Fast part} (\code{frozen\_fast!}), evolving $z$:
\begin{align*}
\partial_t\rho&=-\delta_xU_z-\delta_zW_z,\\
\partial_t\Theta&=-\delta_x\big(\tilde\theta_x(y)\,U_z\big)-\delta_z\big(\tilde\theta_z(y)\,W_z\big),\\
\partial_tU&=-\frac{\Pi_{i,k}-\Pi_{i-1,k}}{\Delta x},\qquad
\partial_tW=-\frac{\Pi_{i,k}-\Pi_{i,k-1}}{\Delta z}-g\,\avg{\rho_z}_z,
\end{align*}
plus the sponge on all rows, with the pressure difference across a face
\[
\text{\code{:gradp}:}\quad \Pi_r-\Pi_l=P(\Theta_{z,r})-P(\Theta_{z,l}),\qquad
\text{\code{:product}:}\quad \Pi_r-\Pi_l=P'_f(y)\,\big(\Theta_{z,r}-\Theta_{z,l}\big).
\]
\paragraph{Slow part} $S(y)=(0,0,-A_U(y),-A_W(y))$, the momentum advection evaluated at $y$ only.

\paragraph{MIS stage.} Stage $i$ integrates
\[
\frac{\mathrm dZ}{\mathrm d\tau}=\sum_jb_{ij}\,F(Y_j,Z)+\sum_j\frac{\gamma_{ij}}{d_i\dt}(Y_j-y^n),\qquad b_{ij}=\frac{\beta_{ij}}{d_i}.
\]
Since $F$ is affine in the coefficients, $\sum_jb_{ij}F(Y_j,Z)=F_{\bar C}(Z)+\sum_jb_{ij}S(Y_j)$ with the weighted
coefficients $\bar C=\sum_jb_{ij}C(Y_j)$: one fast operator per stage
(\code{frozen\_mis\_step!}).

\section{What is frozen, and what must not be}
\begin{center}\small
\begin{tabular}{lll}
\toprule
term & frozen from $y$ & live from $z$\\
\midrule
mass flux & -- & $U$, $W$\\
$\Theta$ flux & $\tilde\theta_x$, $\tilde\theta_z$, upwind direction & $U$, $W$\\
pressure, scheme \code{:gradp} & -- & $P(\Theta)$\\
pressure, scheme \code{:product} & face means $P'_f$ of $\gamma p/\Theta$ & $\jump\Theta_f$\\
gravity & -- & $\avg{\rho}_z$\\
momentum advection (slow) & everything & --\\
\bottomrule
\end{tabular}
\end{center}
Rule of thumb: freeze \emph{coefficients} of the fast operator, never a part of its stiff response. A
third option, $p=(p/\Theta)(y)\,\Theta_z$, was tried and removed: since $\partial p/\partial\Theta=\gamma\,p/\Theta$, it puts
$(\gamma-1)/\gamma$ of the acoustic stiffness into the frozen argument, and the split is then stable only for
$c\,\dt/\Delta z\lesssim1$.

\section{Properties}
\begin{enumerate}
\item $F(q,q)=N(q)$ exactly in both pressure modes (test: relative difference $\lesssim10^{-10}$).
\item For the scheme \code{:product} the fast part is linear in $z$: $F(y,z)=L(y)z+S(y)$, with the coefficients
  $\tilde\theta$, $P'_f$ of $y$; for \code{:gradp} it keeps the cellwise $P(\Theta_z)$.
\item Mass and $\Theta$ are conserved (flux form), the momentum equations are in flux form as well.
\item The $\Theta$ flux of the fast part is the live momentum times the frozen $\theta$,
  $\tilde\theta(y)\,U_z$ and $\tilde\theta(y)\,W_z$; it is linear in $z$, no velocity is formed, and $\tilde\theta(y)$ is a
  coefficient like $P'_f(y)$. The slow part contains no $\Theta$ terms. See Section~\ref{sec:mis}.
\end{enumerate}

\section{Fast integrators}
\paragraph{Explicit.} Forward--backward (Verlet) in the thermodynamic rows $(\rho,\Theta)$ and the momentum rows,
substep limited by the vertical sound speed, $\delta\tau\lesssim\Delta z/c$.
\paragraph{HEVI.} Strang splitting: half a horizontal Verlet step, a vertically implicit step, half a
horizontal step. The vertical part (vertical fluxes of $\rho$, $\Theta$, vertical pressure difference, gravity)
is advanced by the linearly implicit Crank--Nicolson step
$z\leftarrow z+(I-\tfrac{\delta\tau}2J_v)^{-1}\delta\tau\,G_v(z)$. With the pressure difference across the $z$-face $k$
written as $\alpha_k\Theta_k-\beta_k\Theta_{k-1}$, the elimination of $\rho$ and $\Theta$ (as in WRF) gives per column the
tridiagonal system for $W$ ($c=\delta\tau/2$, $r=(c/\Delta z)^2$, $\theta_k=\tilde\theta_z$ at the face $k$):
\begin{multline*}
W_{k-1}\Big[-r\beta_k\theta_{k-1}+\frac{c^2g}{2\Delta z}\Big]+W_k\big[1+r\theta_k(\alpha_k+\beta_k)\big]
+W_{k+1}\Big[-r\alpha_k\theta_{k+1}-\frac{c^2g}{2\Delta z}\Big]\\
=R^W_k-\frac c{\Delta z}\big(\alpha_kR^\Theta_k-\beta_kR^\Theta_{k-1}\big)-\frac{c\,g}2\big(R^\rho_{k-1}+R^\rho_k\big),
\end{multline*}
then $\rho$ and $\Theta$ by back substitution. \code{:product}: $\alpha_k=\beta_k=P'_{k-1/2}(\bar y)$, the vertical
part is linear and the step is exactly Crank--Nicolson. \code{:gradp}: $\alpha_k=a_k$, $\beta_k=a_{k-1}$ with
$a=\gamma p/\Theta$ at the start of the stage (an approximate Jacobian). The factorization is computed once per
stage and step size (\code{frozen\_stage\_begin!}, \code{frozen\_vertical\_step!}).

\section{Well-balanced variant (not implemented)}\label{sec:wb}
At a hydrostatic state the $W$ equation leaves the residual
$-\big(\jump p_{k-1/2}+g\Delta z\,\avg\rho_{k-1/2}\big)/\Delta z$, which vanishes only up to truncation
error with the arithmetic mean $\avg\rho$. Following Artiano, Knoth, Spichtinger, Ranocha (arXiv:2509.10311),
the gravity term is the two-point term $\bar\rho\,\jump\phi$ ($\phi=gz$, $\jump\phi=g\Delta z$) with a mean chosen
such that $\bar\rho\,\jump\phi=-\jump p$ for the hydrostatic state:
\[
\text{isothermal: }\bar\rho=\avg\rho_{\log}=\frac{\jump\rho}{\jump{\log\rho}},\qquad
\text{constant }\theta:\ \bar\rho=\avg\rho_\gamma=\frac{\gamma-1}{\gamma}\frac{\jump{\rho^\gamma}}{\jump{\rho^{\gamma-1}}}.
\]
(Isothermal: $\jump{\log\rho}=-\jump\phi/(RT)$, so $\avg\rho_{\log}\jump\phi=-RT\jump\rho=-\jump p$; constant $\theta$:
$\tfrac{\gamma}{\gamma-1}K\theta^\gamma\rho^{\gamma-1}+\phi=\mathrm{const}$, so $\avg\rho_\gamma\jump\phi=-K\theta^\gamma\jump{\rho^\gamma}=-\jump p$.)
\[
\partial_tW_{i,k-1/2}=-A_W-\frac{\jump p_{k-1/2}+\bar\rho_{k-1/2}\,g\Delta z}{\Delta z}
\]
is then exactly balanced for these atmospheres, without any background state. The $\Theta$ fluxes vanish at
rest ($W=0$), the horizontal pressure differences vanish for horizontally homogeneous states.

\paragraph{Product scheme.} With the product form the hydrostatic balance reads
$P'_{k-1/2}\jump\Theta+\bar\rho\,g\Delta z=0$. For constant $\theta$ ($\Theta=\theta\rho$, $P'=\gamma K\theta^\gamma\rho^{\gamma-1}$,
$\jump{\rho^{\gamma-1}}=-\tfrac{\gamma-1}{\gamma K\theta^\gamma}\,g\Delta z$) the balancing mean is
\[
\bar\rho=(\gamma-1)\,\frac{\avg{\rho^{\gamma-1}}\,\jump\rho}{\jump{\rho^{\gamma-1}}}\quad(\to\rho\ \text{for equal values}).
\]
\paragraph{Frozen form.} Freeze, from $y$, the face derivative $P'_{k-1/2}$ and the ratio
$R_{k-1/2}=\bar\rho_{k-1/2}/\avg\rho_{k-1/2}$, and take $\jump{\Theta_z}$ and $\avg{\rho_z}$ live:
\[
\partial_tW=-\frac{P'_{k-1/2}(y)\,\jump{\Theta_z}_{k-1/2}+R_{k-1/2}(y)\,\avg{\rho_z}_{k-1/2}\,g\Delta z}{\Delta z}.
\]
The fast part stays linear, $F(q,q)=N(q)$, and the column system above only changes its gravity coefficients
($g\to R_{k-1/2}\,g$ in the two $c^2g/(2\Delta z)$ entries and in the $R^\rho$ terms). For a general stratification
(constant Brunt--V\"ais\"al\"a frequency in the gravity-wave tests) neither mean is exact; $\avg\rho_\gamma$ is the
natural default for the troposphere.

\section{Stability of MIS methods for this split}\label{sec:mis}
For a linear problem the frozen split with exact fast stages is MIS with $L=\partial F/\partial z$ and
$N-L=\partial F/\partial y$. Linearized about a hydrostatic state with uniform wind $\bar u$, both operators are
invariant under shifts of the columns, and one horizontal Fourier mode gives a block over the whole column
(\code{MethodDesign/app\_stability.jl}). On the gravity-wave grid ($\Delta x=2.5$\,km, $n_z=16$) the methods
that are stable on the one-dimensional acoustic-advection test of Knoth and Wensch grow by 5--25\,\% per step
already from $C_{\rm adv}=\bar u\dt/\Delta x\approx0.4$. In the frozen split the advection of $\theta$ enters only through the frozen coefficient $\tilde\theta(y)$ of the $\Theta$ flux (Section~5, item 4); this structure is not present in that test.
Methods designed on the column blocks (\code{MethodDesign/search\_app.jl}):
\begin{center}\footnotesize
\begin{tabular}{lcc}
\toprule
method & frozen order & $\max(\rho-1)$, $C_{\rm adv}\le1.6$, $\bar u=10/20/40$ m/s\\
\midrule
MIS54 (Bauer \& Knoth) & 2 & 0.16 / 0.18 / 0.10\\
MISF53 (MIS54 projected onto frozen order 3) & 3 & 0.14 / 0.26 / 0.26\\
MIS\_FV42 (designed, from MIS4) & 2 & $7\cdot10^{-4}$ / $7\cdot10^{-3}$ / $7\cdot10^{-4}$\\
MIS\_FV53 (designed, from MISF53) & 3 & $\le6\cdot10^{-4}$\\
\bottomrule
\end{tabular}
\end{center}
Nonlinear gravity wave (300\,km channel, $U=20$\,m/s, \code{:hevi}): at $T=3000$\,s all are accurate
(MIS\_FV53: wave error $1.2\cdot10^{-4}$\,K at $\dt=130$\,s, $1.8\cdot10^{-4}$\,K at 200\,s; MIS54: $1.6\cdot10^{-4}$,
$1.6\cdot10^{-4}$); in a long run (36000\,s at $\dt=200$\,s, $C_{\rm adv}=1.6$) only MIS\_FV53 survives. The
frozen order (the bicoloured B-series conditions, \code{MethodDesign/BiColourMIS.jl}) is 2 for MIS54 and 3
for MISF53 and MIS\_FV53.

\section{Open points}
\begin{itemize}
\item The split is implemented on flat grids only; terrain (Gal-Chen metric terms in the pressure gradient) and
  cut cells need the frozen coefficients of their metric and aperture terms.
\item The well-balanced gravity mean of Section~\ref{sec:wb} is not implemented.
\item Consistent advection: freezing the transported $\theta$ splits the $\Theta$ advection between the fast and
  the slow part; a split that keeps the advection of $\rho$ and $\Theta$ together (e.g.\ also freezing the
  advective part of the mass flux) would remove the coupling that limits MIS stability.
\item The DG counterpart: \code{docs/frozen\_dg/frozen\_dg.tex}.
\end{itemize}
\end{document}
